%%%%%%%% ICML 2026 Workshop Submission (short / 4-page track) %%%%%%%
%
% Compatible with both AI4Science 2026 (4-8 pages main + unlimited
% refs/appendix) and PhilML 2026 (4 content pages + unlimited refs
% + supplementary).
%
% Uses the official ICML 2026 style; double-blind by default
% (icml2026 strips author info unless [accepted] is set).
% Style files (icml2026.sty / .bst / algorithm.sty / algorithmic.sty
% / fancyhdr.sty) live alongside this .tex file; bibliography at
% pivot_ab_references.bib.
%
%%%%%%%%%%%%%%%%%%%%%%%%%%%%%%%%%%%%%%%%%%%%%%%%%%%%%%%%%%%%%%%%%%%%%

\documentclass{article}

\usepackage{microtype}
\usepackage{graphicx}
\usepackage{subcaption}
\usepackage{booktabs}
\usepackage{multirow}
\usepackage{hyperref}

% Use the following line for the initial blind version submitted for review:
\usepackage{icml2026}

% For preprint, use:
% \usepackage[preprint]{icml2026}
%
% If accepted, instead use the following line for the camera-ready submission:
% \usepackage[accepted]{icml2026}

\usepackage{amsmath}
\usepackage{amssymb}
\usepackage{mathtools}

\usepackage[capitalize,noabbrev]{cleveref}

% Running title for ICML page-headers (under 40 chars).
\icmltitlerunning{Supplementary Material}

\begin{document}

\twocolumn[
  \icmltitle{Supplementary Material for: \\
             Counterfactual Audits Must Precede Conformal Trust for Multimodal Medical Claim Verifiers}

  % Anonymous author block. The icml2026 style file auto-strips this
  % in blind mode (i.e., unless [accepted] is set).
  \begin{icmlauthorlist}
    \icmlauthor{Anonymous Authors}{anon}
  \end{icmlauthorlist}
  \icmlaffiliation{anon}{Anonymous Institution}

  \icmlcorrespondingauthor{Anonymous Authors}{anon@example.org}

  \icmlkeywords{evidence-blindness, multimodal medical verification,
                radiology report generation, counterfactual diagnostic,
                hypothesis-only filter, conformal selection,
                trustworthy AI}

  \vskip 0.3in
]

\printAffiliationsAndNotice{}


% Supplementary material companion document.
% Contains the appendix sections of the main paper, with section
% numbering reset to A, B, C, ... via \appendix.

\appendix
\onecolumn

\section{Threshold Sensitivity for IMG and ESG}
\label{app:thresholds}

The body of the paper treats a verifier as evidence-blind if
$\mathrm{IMG} < 5$\,pp or $\mathrm{ESG} < 5$\,pp. We chose the $5$\,pp
threshold by sweeping $\{3, 5, 7, 10\}$\,pp on the validation split and
selecting the operating point that maximally separated the v5.0 / v5.1
baselines (whose IMG is $\sim 2$\,pp) from the v5.2 / v5.3 / v6.0
configurations (whose IMG is $\geq 60$\,pp). The qualitative finding
(``the highest-validation-accuracy configuration is evidence-blind'') is
robust across all four thresholds tested. ESG is informative on this
benchmark only as a coarse pass/fail signal because every trained
configuration sits in the $[18.56, 21.13]$\,pp band; this is consistent
with the dataset-dependence of partial-input gaps reported by
\citet{poliak2018hypothesis}.

\section{Training Hyperparameters}
\label{app:training}

\begin{table*}[t]
\centering
\small
\caption{Training hyperparameters shared across all reported
configurations. Per-configuration loss-weight settings are derived from
the Methods section of the main paper; full per-configuration YAML is in the
released codebase.}
\begin{tabular}{ll}
\toprule
Image backbone & BiomedCLIP ViT-B/16 (top 4 of 12 blocks trainable) \\
Text backbone & RoBERTa-large (top 8 of 24 layers trainable) \\
Shared dim / fusion layers & $d = 768$, $4$ cross-modal layers, $12$ heads \\
Verdict head & 2-way softmax on a learned $[\mathrm{VERDICT}]$ token \\
Support-score head & sigmoid scalar for conformal calibration \\
Grounding head & $14 \times 14$ patch sigmoid map \\
Optimizer & AdamW \\
LR (encoders / heads) & $1\times 10^{-5}$ / $5\times 10^{-5}$ \\
Weight decay & $0.01$ \\
Warm-up / schedule & $500$-step linear warm-up, cosine decay to zero \\
Batch size / grad accum & $32$ / $2$ \\
Precision & bf16 \\
HO-filter activation threshold & $0.7$ on $p_{\text{text}}^{\text{correct}}$ \\
HO-filter $\mathcal{L}_{\mathrm{cls}}$ downweight & $0.2\times$ \\
$\mathcal{L}_{\mathrm{consist}}$ margin & $0.2$ \\
$\mathcal{L}_{\mathrm{contrast}}$ margin & $0.2$ \\
$\mathcal{L}_{\mathrm{uncert}}$ MC-dropout rate / samples & $0.2$ / $5$ \\
Random seed & $42$ (global) \\
\bottomrule
\end{tabular}
\end{table*}

The full composite loss is
\begin{equation*}
\mathcal{L} = \lambda_{\mathrm{cls}}\mathcal{L}_{\mathrm{cls}}
              + \lambda_{\mathrm{ground}}\mathcal{L}_{\mathrm{ground}}
              + \lambda_{\mathrm{consist}}\mathcal{L}_{\mathrm{consist}}
              + \lambda_{\mathrm{contrast}}\mathcal{L}_{\mathrm{contrast}}
              + \lambda_{\mathrm{uncert}}\mathcal{L}_{\mathrm{uncert}}
\end{equation*}
with $\lambda_{\mathrm{cls}} = 1.0$, $\lambda_{\mathrm{ground}} = 0.5$,
$\lambda_{\mathrm{consist}} = 0.3$, $\lambda_{\mathrm{contrast}} = 0.3$,
$\lambda_{\mathrm{uncert}} = 0.1$.

\section{Loss-Drop Ablations}
\label{app:ablations}

We hold the v5.3 configuration fixed and remove one loss term at a time
to identify the load-bearing components. Table~\ref{tab:ablations}
reports IMG at the resulting test-time test split.

\begin{table*}[t]
\centering
\small
\caption{Single-component drop ablations on the v5.3 configuration. IMG
collapses below the $5$\,pp evidence-blind threshold when either
$\mathcal{L}_{\mathrm{consist}}$ or the HO filter is removed; removing
the grounding, contrastive, or uncertainty terms each changes IMG by
$< 5$\,pp.}
\label{tab:ablations}
\begin{tabular}{lr}
\toprule
\textbf{Variant of v5.3} & \textbf{IMG (pp)} \\
\midrule
Full v5.3 & $69.24$ \\
$-\mathcal{L}_{\mathrm{consist}}$ (\emph{load-bearing}) & $\phantom{0}2.13$ \\
$-$ HO filter (\emph{load-bearing}) & $\phantom{0}5.85$ \\
$-\mathcal{L}_{\mathrm{ground}}$ & $66.10$ \\
$-\mathcal{L}_{\mathrm{contrast}}$ & $63.28$ \\
$-\mathcal{L}_{\mathrm{uncert}}$ & $73.40$ \\
\bottomrule
\end{tabular}
\end{table*}

The two interventions act on disjoint failure modes: the consistency
loss penalises verdict-invariance under image masking (a per-example
signal), and the HO filter downweights training rows the text-only model
already solves (a row-reweighting signal). Both are required.

\section{Benchmark Cohort and Silver-Label Protocol}
\label{app:benchmark}

\paragraph{Training cohort.} \textsc{GroundBench-v6}
aggregates three public sources: OpenI \citep{demner2016openi} (US
frontal radiographs with paired CheXbert-format reports), ChestX-Det10
\citep{liu2020chestxdet} (US radiographs with pixel annotations over
ten finding categories), and PadChest-GR \citep{castro2024padchestgr}
(bilingual Spanish/English studies with per-sentence radiologist
bounding boxes). Patient-stratified splits at $70/10/10/10$ via
MD5-hash bucketing on the patient identifier yield $28{,}361$ training,
$4{,}046$ validation, $4{,}046$ calibration, and $4{,}132$ test claims
after filtering for resolved \textsc{sup}/\textsc{con} labels.

\paragraph{Real-RRG silver labels.} For real-hallucination evaluation we
generate findings-section reports from $500$ held-out OpenI studies
using MAIRA-2 \citep{bannur2024maira2}, CheXagent-2-3B
\citep{chen2024chexagent}, and MedGemma-4B-IT \citep{google2024medgemma}.
Each generated report is decomposed into atomic claims; the three
generators yield $2{,}707$ atomic claims after decomposition. We
silver-label each claim with a three-LLM ensemble (a MIMIC-fine-tuned
grader plus two MIMIC-free graders for leakage decoupling
\citep{ostmeier2024green}) and report two-of-three majority labels.
$1{,}599$ claims resolve to \textsc{sup}, $335$ to \textsc{con}, and
$773$ to \textsc{uncertain}.

\section{Field Audit: Per-Detector Numbers}
\label{app:fieldaudit}

\begin{table*}[t]
\centering
\small
\caption{Zero-shot detectors on the $500$-claim real-RRG subset. RadFlag
is a hallucination flagger that emits a binary stability score and is
not input-conditioned; IMG/ESG do not apply (\textsc{n/a}).}
\label{tab:baselines}
\begin{tabular}{lcccc}
\toprule
\textbf{Detector} & \textbf{Acc} & \textbf{IMG (pp)} & \textbf{ESG (pp)} & \textbf{Blind?} \\
\midrule
BiomedCLIP (zero-shot)              & $0.538$ & $\phantom{0}8.0$  & $\phantom{0}0.0$  & yes \\
MedGemma-4B-IT (as verifier)        & $0.758$ & $14.0$ & $\phantom{0}2.0$  & yes \\
MAIRA-2 (as verifier)               & $0.498$ & $32.2$ & $\phantom{0}6.6$  & no  \\
RadFlag (MAIRA-2, $K{=}6$)          & $0.664$ & \textsc{n/a}  & \textsc{n/a}  & \textsc{n/a} \\
Ours (v6.0-retrain)                 & $0.911$ & $70.21$ & $20.01$ & no  \\
\bottomrule
\end{tabular}
\end{table*}

The combination of axes matters: a verifier can fail one axis without
failing the other, and aggregate accuracy alone does not predict the
joint outcome.

\section{Inverted Conformal Selection Procedure}
\label{app:conformal}

We wrap each trained verifier in an inverted conformal Benjamini--
Hochberg selection procedure calibrated against the support-score
distribution of \textsc{con} claims on a held-out $4{,}046$-claim
calibration split. The procedure follows the conformal-factuality
template of \citet{mohri2024conformal} and the FDR-controlled selection
template of \citet{gui2024conformal}. At target $\alpha = 0.10$, only
v5.3-contrast produces a non-empty trusted set: $n_{\mathrm{green}} =
985$ claims at achieved $\mathrm{FDR} = 0.91\%$ and power $39.6\%$. The
v5.0, v5.1, and v5.2 configurations collapse to $n_{\mathrm{green}} =
0$.

The mechanism is joint-distribution-driven rather than marginal:
the support-score distributions of v5.0 and v5.1 are not less sharp
than v5.3 on the marginal axis, but they fail to separate \textsc{sup}
from \textsc{con} along the score-label joint, so the calibration set's
\textsc{con} scores cover the test set densely and Benjamini--Hochberg
finds no cutoff at $\alpha = 0.10$. This is the empirical statement of
the discussion-section claim that ``conformal selection on an
evidence-blind verifier has nothing to certify.''

\section{Reproducibility Checklist}
\label{app:repro}

\textbf{Met:} (i) Training data sources are public (OpenI, ChestX-Det10,
PadChest-GR); MIMIC-CXR is not used. (ii) Patient-stratified
$70/10/10/10$ splits via MD5-hash bucketing on patient identifier. (iii)
Hyperparameters in Appendix~\ref{app:training}; per-configuration YAML
in the released codebase. (iv) Modal H100 80\,GB; bf16 mixed precision.

\textbf{Deferred to extended version:} model and code release URLs
(anonymized for double-blind review, filled at camera-ready); multi-seed
replication; PadChest-GR claim-extraction for full three-site training.

\end{document}
